\documentclass[11pt]{article}
\usepackage[margin=1in]{geometry}
\usepackage[T1]{fontenc}
\usepackage[utf8]{inputenc}
\usepackage{amsmath,amssymb}
\usepackage{xcolor}
\usepackage{microtype}
\usepackage[colorlinks=true,urlcolor=blue!55!black]{hyperref}

\newcommand{\Ces}{\operatorname{Ces}}
\newcommand{\Cop}{\operatorname{Cop}}

\title{The missing Copson--Ces\`aro embedding regime\\was completely characterized in 2022}
\author{Literature-status note for arXiv:1507.07866}
\date{11 August 2026}

\begin{document}
\maketitle

\begin{center}
\fbox{\parbox{0.91\textwidth}{\textbf{Status: literature\_already\_answered.}
The open regime $p_2>q_2$ in the embedding
$\Cop_{p_1,q_1}(u_1,v_1)\hookrightarrow
 \Ces_{p_2,q_2}(u_2,v_2)$ is covered by Gogatishvili, Pick, and
\"Unver, arXiv:2203.00596 / JFA 283 (2022).  Their Theorem~2.1 gives
necessary and sufficient conditions, with two-sided estimates of the optimal
constant, for the equivalent four-weight Hardy--Copson inequality for all
finite positive parameters.  Remark~2.3 explicitly transfers the theorem
back to the stated embedding.}}
\end{center}

\section*{The source problem}

Gogatishvili, Mustafayev, and \"Unver characterize
\[
 \Cop_{p_1,q_1}(u_1,v_1)\hookrightarrow
 \Ces_{p_2,q_2}(u_2,v_2)
\]
for $0<p_1,p_2,q_1,q_2<\infty$ under $p_2\le q_2$.  In the
Introduction they explain that their duality method works only in that
range and state: ``in the case when $p_2>q_2$ the characterization of these
embeddings remain[s] open'' \cite[p.~3]{source}.  Lemma~7.1 already disposes
of $p_1<p_2$: under the paper's nondegeneracy assumptions no such embedding
can hold.  Thus the substantive omitted range has $p_2>q_2$ and
$p_2\le p_1$.

\section*{The later theorem and the exact bridge}

The later paper studies the same inequality and makes the substitutions
\begin{equation}\label{eq:substitution}
 r=\frac{p_2}{p_1},\qquad q=\frac{q_2}{p_1},\qquad
 p=\frac{q_1}{p_1},
 \qquad
 u=u_2^{q_2},\quad v=v_1^{-p_2}v_2^{p_2},\quad w=u_1^{q_1}.
\end{equation}
After replacing $f$ by $(fv_1)^{p_1}$ and taking the $p_1$-th power, the
embedding inequality is equivalent to
\begin{equation}\label{eq:reduced}
 \left(\int_0^\infty
   \left(\int_0^t f(s)^r v(s)\,ds\right)^{q/r}u(t)\,dt
 \right)^{1/q}
 \le C
 \left(\int_0^\infty
   \left(\int_t^\infty f(s)\,ds\right)^p w(t)\,dt
 \right)^{1/p},
\end{equation}
with $C=c^{p_1}$ for the two optimal constants \cite[pp.~5--6]{answer}.

Theorem~2.1 of the later paper characterizes \eqref{eq:reduced} for every
$0<r\le1$ and every $0<p,q<\infty$, in seven exhaustive parameter regions,
and estimates its best constant in each region \cite[pp.~6--8]{answer}.
The old missing hypothesis $p_2>q_2$ is exactly $q<r$.  Within $r\le1$,
its possibilities are exhausted as follows:
\[
 p\le q<r \quad\text{(case (ii))},\qquad
 q<p\le r \quad\text{(case (v))},\qquad
 q<r<p \quad\text{(case (vi))}.
\]
Remark~2.3 then states explicitly that the best constant in the original
Copson-to-Ces\`aro embedding is obtained from Theorem~2.1 by the substitutions
\eqref{eq:substitution} \cite[p.~8]{answer}.  This is a full answer, not merely
a sufficient condition or an answer for special weights.

\section*{Independent confirmation and scope}

The authors' 2024 survey identifies the old $p_2\le q_2$ results and says
that the complete characterization of the same embedding, ``without any
restrictions on parameters or weights,'' was given by the 2022 JFA paper
\cite[p.~4]{survey}.  The conclusion here is therefore a literature-status
identification; no new proof of the seven-case criterion is claimed.

The statement concerns finite positive exponents, exactly the parameter
range of arXiv:1507.07866.  The endpoint problem with infinite exponents is
not needed to answer that source question.

\sloppy
\begin{thebibliography}{3}
\raggedright
\bibitem{source}
A.~Gogatishvili, R.~Mustafayev, and T.~\"Unver,
\emph{Embeddings between weighted Copson and Ces\`aro function spaces},
arXiv:1507.07866; Czechoslovak Math. J. \textbf{67} (2017), 1105--1132.

\bibitem{answer}
A.~Gogatishvili, L.~Pick, and T.~\"Unver,
\emph{Weighted inequalities involving Hardy and Copson operators},
arXiv:2203.00596; J. Funct. Anal. \textbf{283} (2022), Paper 109719.

\bibitem{survey}
A.~Gogatishvili and T.~\"Unver,
\emph{On weighted Ces\`aro function spaces}, arXiv:2410.22301v2 (2025).
\end{thebibliography}

\end{document}
